% ================================================================
%  ДҮГНЭЛТ
% ================================================================
\chapter*{ДҮГНЭЛТ}
\addcontentsline{toc}{chapter}{Дүгнэлт}

Энэхүү дипломын ажлын хүрээнд ``Рестораны ширээ захиалга, цэс удирдлагын мобайл
системийн загварчлал, хэрэгжилт'' сэдвийн дагуу QR кодоор суурилсан бүрэн
автоматжсан мобайл системийг амжилттай загварчилж хэрэгжүүлсэн.

Тавьсан зорилтуудыг дараах байдлаар биелүүлсэн:

\begin{enumerate}[leftmargin=*]
  \item \textbf{Онолын үндэслэл тогтоосон:} QR код (ISO/IEC 18004:2015), WebSocket
    (RFC 6455), REST API, JWT (RFC 7519) зэрэг технологийн онолын судалгааг хийж,
    Toast POS, Lightspeed, Square, Ajisen Ramen зэрэг ижил төстэй системүүдтэй
    харьцуулсан дүн шинжилгээ хийсэн;

  \item \textbf{Хэрэглэгч болон техникийн шаардлагыг тодорхойлсон:} 4 хэрэглэгчийн
    үүрэг (Зочин, Кассир, Менежер, Гал тогоо), 12 функционал шаардлага (MoSCoW
    чухлын зэрэглэлтэй), 5 ангилалд бүлэглэсэн функционал бус шаардлагыг ISO/IEC
    25010 стандартад нийцүүлэн тодорхойлсон;

  \item \textbf{Системийн загварыг боловсруулсан:} QR кодоос захиалга ажилтанд
    хүртэлх бүрэн урсгалын загварыг ER диаграм (9 хүснэгт), компонент, хэрэгжүүлэлт,
    захиалгын state machine, 3 дарааллын диаграм, swimlane бүхий үйл ажиллагааны
    диаграм, 10 юзкейс тодорхойлолт бүхий UML юзкейс диаграмаар тодорхойлсон;

  \item \textbf{4 давхаргат архитектурыг хэрэгжүүлсэн:} React 18 + TypeScript
    (Presentation), Node.js/Express.js (Application), JWT + Drizzle ORM (Business
    Logic), PostgreSQL 16 (Data) давхаргуудаар бүрдсэн системийг pnpm workspace
    монорепо бүтцээр хөгжүүлж, 30 REST API endpoint, 9 Socket.IO үйл явдал
    хэрэгжүүлсэн;

  \item \textbf{Туршилт хийж, үр дүнд дүн шинжилгээ хийсэн:}
    \begin{itemize}
      \item k6 ачааллын туршилт: API хариу дунджаар 120 мс ($<$ 200 мс шаардлага),
        WebSocket дамжуулалт 45 мс ($<$ 100 мс шаардлага);
      \item SUS хэрэглэгчийн үнэлгээ: 84.2/100 оноо (``Маш сайн'' зэрэглэл, A grade);
      \item Уламжлалт аргатай харьцуулалт: захиалга өгөх хугацаа 62\%, алдааны хувь
        94\%-иар буурсан.
    \end{itemize}
\end{enumerate}

\section*{Системийн давуу тал}

\begin{itemize}
  \item \textbf{Хямд өртөг}: Сарын зардал $\sim$10--25 ам.доллар (гадаадын системүүдийн
    60--165 ам.доллартай харьцуулахад);
  \item \textbf{Монгол хэлний бүрэн дэмжлэг}: Цэс, интерфейс, мэдэгдлүүд
    бүгд Монгол хэлээр;
  \item \textbf{Нээлттэй эх код}: Рестораны шаардлагад нийцүүлэн тохируулах боломжтой;
  \item \textbf{Бодит цагийн харилцаа}: WebSocket/Socket.IO-р захиалга шуурхай дамжина;
  \item \textbf{Суулгахад хялбар}: QR код уншуулахад л болох тул хэрэглэгчдэд тусгай
    програм суулгах шаардлагагүй.
\end{itemize}

\section*{Цаашид хөгжүүлэх боломж}

\begin{itemize}
  \item QPay, SocialPay, Monpay зэрэг Монголын цахим төлбөрийн системтэй холбох;
  \item Хиймэл оюун ухааны (AI) санал болгох систем нэвтрүүлэх;
  \item Олон салбар рестораны удирдлагын модуль нэмэх;
  \item AR (Augmented Reality) технологиор хоолны 3D загвар үзүүлэх;
  \item Гар утасны native аппликейшн (React Native) хөгжүүлэх;
  \item Нөөцийн оновчлол (capacity planning) модуль нэмэх.
\end{itemize}

\section*{Дүгнэлтийн товчлол}

Энэхүү дипломын ажлаар Монгол Улсын рестораны салбарт тохирсон, хямд өртөгтэй,
нэвтрүүлэхэд хялбар, Монгол хэлийг бүрэн дэмждэг QR кодод суурилсан мобайл
захиалгын системийг амжилттай хөгжүүлж, туршилтаар баталгаажуулсан. Тус систем нь
рестораны үйлчилгээний чанарыг дээшлүүлж, зардал хэмнэж, хэрэглэгчийн сэтгэл
ханамжийг нэмэгдүүлэх бодит хэрэгсэл болохыг харуулсан.
